% =========================================================================
% Chapter 7: Conclusion
% =========================================================================

\chapter{Conclusion}
\label{ch:conclusion}

Whole plants of \textit{Dipsacus inermis} were obtained from the Department of Botany, University of Kashmir (34.1280$^{\circ}$\,N, 74.8365$^{\circ}$\,E), identified, washed, shade-dried, and powdered. An aqueous extract was prepared and filtered. Qualitative tests were positive for the major secondary-metabolite classes screened. Clinical pathogenic stocks of \textit{Escherichia coli}, \textit{Pseudomonas aeruginosa}, and \textit{Klebsiella pneumoniae} from the Advanced Research Laboratory, Department of Zoology, were tested by disc diffusion at C1 (\qty{50}{\milli\gram\per\milli\litre}) and C2 (\qty{100}{\milli\gram\per\milli\litre}). Zones were measured inclusive of the discs. Distilled-water controls were zero. Gentamicin 5 discs produced marked inhibition on every plate. Extract discs produced clear zones against all three organisms; the response was concentration-dependent for \textit{E.\ coli} and \textit{K.\ pneumoniae}.

The three stated objectives were met: collection and preparation of the aqueous extract; preliminary phytochemical profiling; and in vitro antibacterial screening against the three clinical pathogenic strains. The work does not show clinical efficacy, a minimum inhibitory concentration, or a characterised lead compound \citep{vanvuuren2019,balouiri2016}.

Useful next measurements would be a deposited voucher, strain accession numbers, Mueller--Hinton MIC, instrumental chemistry of the same aqueous batch, and a photographed \textit{P.\ aeruginosa} inhibition series.
